\subsection*{导体边界上的局部规律}

静电平衡时，导体表面电场沿法线，但不同位置的场强一般不同。由
$E_n=\sigma/\varepsilon_0$ 可知，表面曲率较大、尖锐程度较高的位置往往需要更高的面电荷密度
来维持整个导体等势，因而附近场强更大。这是尖端放电和避雷针附近容易发生气体电离的基础。
“电荷集中在尖端”不是额外定律，而是等势边界条件与几何共同决定的结果。

在导体表面取一个电荷元 $\sigma\dd S$，它不能对自身施力。表面外侧总场为
$E_{\rm out}=\sigma/\varepsilon_0$，内侧为零；去除该电荷元自身在两侧方向相反的场后，其他电荷
对它的平均场为
\begin{equation}
 E_{\rm others}=\frac{E_{\rm out}+E_{\rm in}}2
 =\frac\sigma{2\varepsilon_0}.
\end{equation}
因此导体表面受到向外的静电压强
\begin{equation}
 \boxed{p=\sigma E_{\rm others}=\frac{\sigma^2}{2\varepsilon_0}
 =\frac12\varepsilon_0E_{\rm out}^2}.
\end{equation}
它与真空电场能量密度相等，可用于求带电球壳的膨胀趋势和电容器极板间的吸力。

\subsection*{空腔、有内电荷与接地}

若闭合导体空腔内放入电荷 $q$，在包围空腔而仍处于导体材料内部的高斯面上 $\vect E=0$，故
空腔内壁的总感应电荷必须为 $-q$。若导体原来总电荷为 $Q$，外表面的总电荷便为 $Q+q$。
这些结论只确定各表面的总量；具体面密度仍由空腔形状、内电荷位置和外部边界条件决定。

内电荷会使空腔中出现电场，却不会破坏导体材料内部的零场。外部场在只给定外表面总电荷和外部
边界以后，也不需要知道空腔中电荷的具体位置。这是静电屏蔽的另一层含义：闭合导体把内外两个
边值问题分开。若导体接地，外表面还可与大地交换任意电荷，外部场甚至可能完全不显露腔内电荷。

接地条件不是“导体电荷为零”，而是导体电势固定为大地电势，通常取零。接地导体可以带有非零
净电荷；孤立中性导体则必须保持总电荷为零，但其电势一般不为零。电像法中若把接地球改为孤立
中性球，原来的单个像电荷只能满足球面等势，不能同时满足总感应电荷为零；还需在球心再加入一个
辅助像电荷，使所有像电荷代数和为零。

\subsection*{唯一性定理的使用步骤}

用唯一性定理解决导体问题时，应依次检查：所研究区域内的真实体电荷是否相同；每个边界上的
电势或法向场条件是否相同；无穷远条件是否相同。只要三者满足，构造解便是该区域的唯一解。
像电荷可以放在研究区域之外，例如导体材料内部，因为那里并不要求构造场等于真实场；若把像
电荷误放入要求求解的真空区域，就改变了 Poisson 方程的电荷源，不能再用唯一性定理证明等价。

接地平面外点电荷是最简单的例子。把平面取为 $z=0$，真实电荷 $Q$ 位于 $z=a$，在 $z=-a$
放置像电荷 $-Q$。上半空间电势
\begin{equation}
 \varphi=\frac1{4\pi\varepsilon_0}
 \left(\frac Q{r_+}-\frac Q{r_-}\right)
\end{equation}
在平面上处处为零，且上半空间只有真实电荷，因此就是所求解。真实电荷受力等于它与像电荷间的
库仑力：
\begin{equation}
 \boxed{F_z=-\frac{Q^2}{16\pi\varepsilon_0a^2}}.
\end{equation}
导体表面感应电荷密度的积分为 $-Q$。这并不意味着一块有限孤立金属板必然感应出总量 $-Q$；
该结论依赖无限接地平面的边界条件。

\subsection*{像电荷能量中的二分之一}

求真实电荷受到的瞬时力时，导体当时已经完成感应，故可把像电荷当作固定场源直接使用，不乘
$1/2$。求系统能量时，感应电荷却是随着真实电荷从无穷远移入而逐渐建立的。以接地平面为例，
像电荷与真实电荷相距 $2a$，若直接写两点电荷相互作用能会得到
$-kQ^2/(2a)$；真实系统能量应为其一半：
\begin{equation}
 U=-\frac{kQ^2}{4a}=-\frac{Q^2}{16\pi\varepsilon_0a}.
\end{equation}
对 $a$ 求负导数恰好恢复上面的吸力。力和能量两种方法给出同一结果，是检查二分之一因子是否
处理正确的可靠办法。

\subsection*{等离子体振荡的物理尺度}

$\omega_p$ 随数密度的平方根增大，随载流子质量的平方根减小。电子比离子轻得多，所以高频响应
通常由电子主导，而离子在一个电子振荡周期内近似静止。当外加电场频率低于或接近等离子体频率
时，电子有足够时间重排并显著屏蔽外场；频率远高于 $\omega_p$ 时，电子来不及跟随，介质的响应
性质会改变。

小位移推导假设电子云整体平移，回复力与位移成正比。若振幅太大、密度明显不均匀、碰撞频繁或
存在磁场，振荡会出现阻尼、色散和不同模式，单一的简谐频率只给出最基本的数量级。尽管如此，
$\omega_p=\sqrt{ne^2/(m\varepsilon_0)}$ 仍是判断等离子体时间尺度的首要参数。
